% ===================== Chapter 13 =====================
\chapter{Closures and Lambdas}
\label{ch:closures-lambdas}

So far functions have been things we \emph{define} and \emph{call}. In this
chapter functions become \emph{values}: things we can store in a variable, pass to
another function, and create on the spot without giving them a name. A small,
unnamed function is called a \textbf{lambda}, and when one captures variables from
around it, it becomes a \textbf{closure}. These ideas unlock a flexible,
expressive style of programming.

\section{Functions as values}

A \textbf{lambda} is a function written as an expression, with no name. Its syntax
is a parameter list, the arrow \code{=>}, and a body. For a one-expression body,
just write the expression:

\begin{salam}
func main:
    double auto = (x i32) => x * 2
    println double(21)
end
\end{salam}

\begin{progout}
42
\end{progout}

Here \code{(x i32) => x * 2} is a function value: it takes an \code{i32} called
\code{x} and produces \code{x * 2}. We store it in the variable \code{double}
(letting \code{auto} infer its type) and then \emph{call} it just like any
function: \code{double(21)}. The lambda has no name of its own the name
\code{double} belongs to the variable, not the function.

\subsection{The type of a function value}

The type of \code{double} is \code{func(i32) i32}: ``a function taking an
\code{i32} and returning an \code{i32}.'' In general a function-value type is
written \code{func(}\textit{argument types}\code{) }\textit{return type}. A
function value taking nothing and returning an \code{i32} has type
\code{func() i32}. You will see these types most often where one function accepts
another as a parameter.

\subsection{Multi-statement bodies}

When a lambda needs more than one statement, give it a block body: a colon after
the arrow, the statements, and \code{end} exactly like a named function, just
without the name:

\begin{salam}
(x i32) => :
    mut y i32 = x * x
    ret y + 1
end
\end{salam}

The single-expression form (\code{=> x * 2}) is shorthand for the common case; the
block form (\code{=> : ... end}) is there when you need it.

\section{Higher-order functions}

A function that takes another function as a parameter or returns one is
called a \textbf{higher-order function}. Because a function value has a type like
\code{func(i32) i32}, you can declare a parameter of that type and call it inside:

\begin{salam}
func apply(f func(i32) i32, x i32) i32:
    ret f(x)
end

func main:
    println apply((x i32) => x + 100, 5)
end
\end{salam}

\begin{progout}
105
\end{progout}

\code{apply} does not know or care what \code{f} does; it just calls \code{f(x)}.
We hand it a lambda created right at the call site, \code{(x i32) => x + 100}, and
\code{apply} runs it on \code{5}. This is the essence of higher-order
programming: behaviour passed around as data. The same \code{apply} could square,
negate, or double its argument, depending only on the lambda you give it.

\section{Closures: capturing the surrounding state}

Here is where function values become truly powerful. A lambda can \emph{use
variables from the scope in which it was created} not just its own parameters.
When it does, it ``closes over'' those variables, and we call it a
\textbf{closure}. The captured values travel with the function, even after the
scope that created them has returned.

The classic example is a function that \emph{builds} another function:

\begin{salam}
func make_adder(n i32) func(i32) i32:
    ret (x i32) => x + n
end

func main:
    add10 auto = make_adder(10)
    add100 auto = make_adder(100)
    println add10(5)
    println add100(5)
end
\end{salam}

\begin{progout}
15
105
\end{progout}

\code{make\_adder} returns a lambda that adds \code{n} but \code{n} is
\code{make\_adder}'s parameter, not the lambda's. The returned closure
\emph{captures} the value of \code{n}. So \code{make\_adder(10)} produces a
function that always adds 10, and \code{make\_adder(100)} produces a separate one
that always adds 100. Each closure remembers its own captured \code{n}, long after
\code{make\_adder} has finished running.

\section{Closures that remember and change state}

Because a closure carries its captured variables with it, it can also \emph{update}
them between calls giving a function that has memory. A counter is the
canonical example:

\begin{salam}
func make_counter() func() i32:
    mut count i32 = 0
    ret () => :
        count = count + 1
        ret count
    end
end

func main:
    next func() i32 = make_counter()
    println next()
    println next()
    println next()
end
\end{salam}

\begin{progout}
1
2
3
\end{progout}

Each call to \code{next()} increments and returns the captured \code{count}. The
variable \code{count} lives on inside the closure even though \code{make\_counter}
returned long ago the closure has kept it alive. Call \code{make\_counter()} a
second time and you get a fresh, independent counter starting again at zero.

\begin{deep}[In Depth: how closures are implemented]
A closure is more than a pointer to code it must also carry the variables it
captured. Salam implements this by allocating a small \emph{environment} on the
heap holding the captured values, and pairing it with the lambda's code. The
function value you pass around is really this code-plus-environment bundle. When
you call it, the captured variables are read from (and written to) that
environment, which is why \code{count} survives between calls and why two counters
do not interfere. The compiler performs this ``lifting'' of the lambda and its
captures automatically; you simply write the closure and it works.
\end{deep}

\section{Why closures matter}

Closures let you parameterise \emph{behaviour}, not just values. A sort routine
can take a closure describing how to compare two elements; a retry helper can take
a closure describing the work to attempt; an event system can take closures to run
when something happens. In each case you hand over a small piece of custom logic,
created exactly where it is needed and carrying whatever surrounding context it
requires. This keeps general machinery (the sorter, the retrier) cleanly separated
from the specific behaviour you plug into it.

\begin{tip}
Lambdas are at their best when they are small a comparison, a transformation, a
quick test. If a lambda grows into a substantial block of logic, give it a real
name with \code{func} instead; a named function is easier to read, reuse, and
discuss. Reach for a lambda when naming it would add noise, not clarity.
\end{tip}

\section{Chapter summary}

\begin{itemize}[leftmargin=1.4em]
  \item A \textbf{lambda} is an unnamed function value:
        \code{(params) => expression}, or \code{(params) => : ... end} for a
        block body.
  \item A function value has a type like \code{func(i32) i32} or \code{func()
        i32}.
  \item A \textbf{higher-order function} takes or returns a function value; call a
        function-typed parameter just like any function.
  \item A \textbf{closure} is a lambda that captures variables from its
        surrounding scope; the captured values live on with the closure.
  \item Closures can carry mutable state (e.g. a counter), and each closure
        instance has its own independent captured variables.
  \item Use lambdas for small, local behaviour; promote anything substantial to a
        named function.
\end{itemize}

\section{Exercises}

\exercise Create a lambda \code{square} that returns its argument times itself,
store it with \code{auto}, and print \code{square(7)}.

\exercise Write a higher-order function \code{apply\_twice(f func(i32) i32, x
i32) i32} that returns \code{f(f(x))}, and test it with a lambda that adds 3.

\exercise Write \code{make\_multiplier(factor i32) func(i32) i32} that returns a
closure multiplying by \code{factor}. Build a ``times 3'' and a ``times 10''
function from it.

\exercise Using \code{make\_counter}, create two independent counters and show by
interleaving their calls that they do not share state.

\exercise Write a higher-order function \code{transform\_all} that takes a
\code{Vector<i32>} and a \code{func(i32) i32}, and returns a new vector with the
function applied to every element. Test it with a doubling lambda.
